\section{Periods, compact subgroups, and the transcendence question}

We use a Weierstrass equation over $\overline\Q$ and its invariant
differential. Its period lattice is $\Lambda=\Z\omega_1+\Z\omega_2$,
with $\Im(\omega_2/\omega_1)>0$. The quasi-periods are defined by
$\zeta(z+\omega_j)=\zeta(z)+\eta_j$. We extend
$\eta:\Lambda\to\C$ real-linearly to $\C$. The discriminant in
this section is the discriminant of this equation, equivalently of
the lattice in the specified differential. The quantity $N$ is
independent of changing this algebraic equation and differential.

\begin{proposition}[sigma expression and the third-kind period]\hypertarget{tag-0017}{}\label{ari:C.7.1}\hypertarget{tag-IA01}{}\label{ari:C.7.1:lift-transport}
\StatementRecord{ari:C.7.1}{proved}{0017}{\SeriesRef{ari:eq:cp-periods}; \citet[Definitions 2.3--2.4 and
Section 6]{bertolin2022}}

Let $P\in E(\overline\Q)$, with $P\ne O$ and $2P\ne O$, and
let $u$ be an elliptic logarithm of $P$. Then
\begin{equation}\hypertarget{tag-003K}{}\label{ari:eq:cp-sigma}
 2\log N(P)=\frac16\log|\Delta|
       +\Re\{\Lambda_1+u(\zeta(u)-\eta(u))\},
 \qquad \Lambda_1=2\log\sigma(u)-u\zeta(u).
\end{equation}
For a fixed elliptic logarithm $u$, changing the branch of
$\log\sigma(u)$ changes the expression inside the real part by an
integer multiple of $4\pi i$. Put
$B(u)=\Lambda_1+u(\zeta(u)-\eta(u))$. Under a lattice translation,
\[
 B(u+\omega)-B(u)\equiv
 u\eta(\omega)-\omega\eta(u)+2\log\chi(\omega)
 \pmod{4\pi i\Z},\qquad \omega\in\Lambda,
\]
where $\chi(\omega)\in\{1,-1\}$ is the sigma semicharacter.
This difference is purely imaginary, so the displayed real part
is independent of both choices.
There is an extension $G_P$ of $E$ by $\mathbb G_m$, defined
over $\overline\Q$ and corresponding, under the principal
polarization, to $P$, for which $\Lambda_1$, modulo a logarithm
of an algebraic number, is the negative of a relative period
of $[\Z\to G_P]$. Changing an algebraic differential of the
third kind adds an algebraic multiple of $u$ and algebraic
endpoint terms.
\end{proposition}


\begin{proof}
\begin{enumerate}[label=\arabic*.,leftmargin=*]
\item The sigma--theta product, or its divisor and automorphy law
together with its derivative at zero, gives
\begin{equation}\hypertarget{tag-003F}{}\label{ari:eq:cp-klein}
 N(P)=|\Delta|^{1/12}|\sigma(u)|
                   \exp\!\left(-\frac12\Re(u\eta(u))\right).
\end{equation}
To check the constant, for $\Lambda=\Z+\tau\Z$ one uses
$\sigma(z)=e^{\eta_1z^2/2}\vartheta_1(z,\tau)/(2\pi\eta(\tau)^3)$,
$\Delta=(2\pi)^{12}\eta(\tau)^{24}$, and
$\eta(z)=\eta_1z-2\pi i\Im z/T$. Substitution gives exactly
the defining Gaussian in $N$. Homogeneity gives the general
lattice formula. Taking logarithms proves \eqref{ari:eq:cp-sigma}.
The automorphy law
$\sigma(u+\omega)=\chi(\omega)
 e^{\eta(\omega)(u+\omega/2)}\sigma(u)$ and real-linearity of
$\eta$ give the stated lift transport. Explicitly, writing
$u=a\omega_1+b\omega_2$ with $a,b\in\R$ and
$\omega=m\omega_1+n\omega_2$ with $m,n\in\Z$, Legendre's
relation gives
$u\eta(\omega)-\omega\eta(u)=2\pi i(bm-an)$.
This also proves that the transported difference is purely imaginary.

\item We make the period identification explicit. Put
\[
 f_u(z)=\frac{\sigma(z+u)}{\sigma(z)\sigma(u)}e^{-\zeta(u)z}.
\]
Its logarithmic differential is
\[
 \frac{df_u}{f_u}
 =\frac12\frac{\wp'(z)-\wp'(u)}{\wp(z)-\wp(u)}\,dz.
\]
This is the pullback of an algebraic differential with residues
$+1$ at $-P$ and $-1$ at $O$.
Here the principal polarization identifies $P$ with
$t_P^*\mathcal O_E(O)\otimes\mathcal O_E(O)^{-1}
=\mathcal O_E((-P)-(O))$; this fixes the sign of the extension class. The apparent $0/0$ at $z=u$
has a removable singularity when $P$ is not $2$-torsion.
The rational cocycle $f_u(z+z')/(f_u(z)f_u(z'))$ defines
the algebraic extension; in its rational trivialization the
exponential map has the form
\[
 \exp_{G_P}(z,v)=(\exp_E z,e^v f_u(z)).
\]
Its periods include
\begin{equation}\hypertarget{tag-003G}{}\label{ari:eq:cp-periods}
 (0,2\pi i),\qquad
 (\omega_j,\zeta(u)\omega_j-\eta_j u),\quad j=1,2.
\end{equation}
These signs follow directly from the sigma automorphy law.
The identity
$\sigma(2u)=-\wp'(u)\sigma(u)^4$ implies
\[
 \log f_u(u)=\Lambda_1+\log(-\wp'(u)).
\]
For an algebraic lift $R=(P,r)$, $r\in\overline\Q^{\times}$,
the second coordinate of $\log_{G_P}R$ is
\begin{equation}\hypertarget{tag-003L}{}\label{ari:eq:cp-v}
 v=\log r-\Lambda_1-\log(-\wp'(u)).
\end{equation}
It is a period of the relative homology defining the indicated
$1$-motive. This also avoids treating an integral with an endpoint
at a residue pole as an unregularized ordinary integral.
The description agrees with the generalized logarithm and
Serre function in \citet[Definitions 2.3--2.4 and
Section 6]{bertolin2022}; the signs in
\eqref{ari:eq:cp-periods} are for the exponential rather than the
quasi-periods of $f_u$ itself.
\end{enumerate}
\end{proof}

\begin{proposition}[the compact-subgroup criterion and the real formula]\hypertarget{tag-001B}{}\label{ari:C.7.2}
\StatementRecord{ari:C.7.2}{proved}{001B}{\SeriesRef{ari:C.7.1}, \SeriesRef{ari:eq:cp-periods}, \SeriesRef{ari:eq:cp-v}, \SeriesRef{ari:eq:cp-chi}; \citet[slides 5--8]{bertrand2017}}

For $P$ as in Statement \ref{ari:C.7.1}, $N(P)$ is algebraic if and only
if the fibre of $G_P$ above $P$ contains an algebraic point of
the maximal compact subgroup $G_P(\C)^c$.
For the family in the stated conventions, if $t>1$ is real, choose
$\tau=iT$, $T>1$, and the zero $\mu=1/2-iy$, $0<y<T/2$.
With $x=2y/T$, one has
\[
 N(P)^4=3^{3/4}(t-1)^{-1/4}
       \left|\frac{\vartheta_1(x,i/T)}{\eta(i/T)}\right|.
\]
\end{proposition}


\begin{proof}
\begin{enumerate}[label=\arabic*.,leftmargin=*]
\item On the analytic universal cover of $G_P$ define the real
homomorphism
\[
 \chi(z,v)=\Re\{v-\zeta(u)z+u\eta(z)\}.
\]
It vanishes on the periods \eqref{ari:eq:cp-periods}, hence descends
to $G_P(\C)$. Its restriction to the torus is $\log|\cdot|$.
The kernel is compact: in the universal cover it is a real
three-dimensional vector space modulo the rank-three period
lattice. Every compact subgroup has zero image in $\R$, so
this kernel is exactly $G_P(\C)^c$.
The identity
$\Re(u\eta(z))=\Re(z\eta(u))$ follows from the Legendre relation;
equivalently $\eta(z)=sz+\pi\bar z/\operatorname{area}(\Lambda)$
for a complex constant $s$. Substituting
\eqref{ari:eq:cp-v} at $z=u$ gives
\begin{equation}\hypertarget{tag-003E}{}\label{ari:eq:cp-chi}
 \chi(R)=\log|r|-\log|\wp'(u)|
            -2\log N(P)+\frac16\log|\Delta|.
\end{equation}
If $r$ is algebraic and $\chi(R)=0$, this equation makes
$N(P)^2$ algebraic. Conversely, if $N(P)$ is algebraic,
choose the positive real algebraic number
$r=|\wp'(u)|\,|\Delta|^{-1/6}N(P)^2$.
Then \eqref{ari:eq:cp-chi} vanishes. This proves both directions
with the chosen rational trivialization of the fibre.
The compact interpretation is the one discussed by
\citet[slides 5--8]{bertrand2017}; here the calculation fixes
its normalization on the diagonal $P=P$.

\item For the real formula, the rectangular elliptic curve has its
two zeros of $\wp$ on $1/2+i\R$; along that edge $\wp$ runs
between the positive root $e_1$ and the negative root $e_2$.
Thus the indicated representative exists. Translate $2\mu$
by $-1$ and apply the modular substitution $\tau\mapsto-1/\tau$.
The resulting elliptic argument is $-2y/T$, which is real.
The Gaussian is therefore one. Evenness of $N$ and the
duplication identity prove the displayed formula. This removes
the Gaussian from this particular expression but makes no
algebraicity assertion about its theta value.
\end{enumerate}
\end{proof}

\subsection{Motivic groups from the point and extension data}

\citet[Construction~3.4, p.~12, Lemma~3.5, p.~13, and
Proposition~4.12, pp.~19--20]{bertolin2022} compute the motivic
group from its abelian point and extension data and the central
Weil-pairing bracket. For the $1$-motives used here, let $B$ be the
smallest abelian subvariety generated by these data, and let $Z(1)$
be the central torus obtained from that bracket and the remaining
algebraic lift. Their dimension formula is
\begin{equation}\hypertarget{tag-003D}{}\label{ari:eq:cp-MTdim}
 \dim G_{\mathrm{mot}}(M)
 =\dim G_{\mathrm{mot}}(A)+2\dim B+\dim Z(1).
\end{equation}
The bracket is described in their Section~6.1, pp.~26--27, and
their corrected Table~2, pp.~28--29, gives unipotent dimension
three in the dependent nondeficient case. The following calculation
determines these groups through the displayed polarization maps;
the arithmetic question for the metric norm is recorded in the
proved implication with its explicitly open hypothesis in \cref{ari:C.7.4}.

\begin{proposition}[the motivic groups needed on the diagonal]\hypertarget{tag-001C}{}\label{ari:C.7.3}
\StatementRecord{ari:C.7.3}{proved}{001C}{\SeriesRef{ari:eq:cp-MTdim}}

Let $P$ be nontorsion and let $R$ be any algebraic lift above
$P$ in $G_P$. Put $M_R=[\Z\to G_P]$, $1\mapsto R$.
Its connected motivic group is, after a choice of Levi subgroup,
the semidirect product of $G_{\mathrm{mot}}(E)$ and the
three-dimensional Heisenberg group attached to $H_1(E)$ and
its principal polarization. Its dimension is $5$ for CM and
$7$ without CM.

For a non-real embedding put
$\mathcal M=M_R\oplus\overline{M_R}$.
If $E$ and $\bar E$ are not isogenous, its dimension is $13$.
If they are isogenous, choose an isogeny
$\phi:\bar E\to E$ of degree $d$ and write
$Q=\phi(\bar P)$. Let $r$ be the dimension over
$F=\End(E)\otimes\Q$ of the span of $P,Q$, modulo torsion.
Then
\[
\begin{array}{c|c|c|c}
\text{case}&\dim B&\dim Z&\dim G_{\mathrm{mot}}(\mathcal M)\\ \hline
r=2,\ E\text{ CM}&2&2&8\\
r=2,\ E\text{ not CM}&2&2&10\\
r=1,\ E\text{ CM}&1&1+e&5+e\\
r=1,\ E\text{ not CM}&1&1+e&7+e.
\end{array}
\]
Here $e$ is either zero or one: it is zero precisely when the
algebraic lift in the one-dimensional quotient of the two
central tori by the bracket torus is torsion. This criterion is
independent of the choices made modulo isogeny. In the real
case we use the single motive $M_R$; no extra quotient torus
is necessary.
\end{proposition}


\begin{proof}
\begin{enumerate}[label=\arabic*.,leftmargin=*]
\item All curves, polarizations, points, and lifts are algebraic over
$\overline\Q$, as required by Construction 3.4 and Lemma 3.5.
For $M_R$ the smallest abelian subvariety containing $(P,P)$ is
the diagonal $B=E$, since $P$ is nontorsion. If $\psi$ is the
polarization pairing, the bracket on this diagonal is
\[
 [(x,x),(y,y)]=\psi(x,y)-\psi(y,x)=2\psi(x,y).
\]
It is nonzero and spans $\Q(1)$. Hence $Z'=Z=\mathbb G_m$;
there is no deficient or Ribet exception, including for CM.
The pure group is $\mathrm{GL}_2$ without CM and
$\operatorname{Res}_{F/\Q}\mathbb G_m$ for the imaginary
quadratic CM field. Their dimensions are four and two.
This proves the single-motive assertion, including its Lie
bracket and Levi action: the action on the centre is the
polarization multiplier.

\item For the direct sum only the two diagonal central coordinates
occur. If the two elliptic curves are not isogenous, they
cannot be CM conjugates. The pure group is
$\{(g_1,g_2)\in\mathrm{GL}_2^2:\det g_1=\det g_2\}$,
of dimension seven. The point data generate $E\times\bar E$,
and the two brackets are independent, giving $7+4+2=13$.

\item If the curves are isogenous, transport the second polarized
factor through $\phi$. Its polarization pairing becomes
$\psi/d$. The smallest abelian subvariety has dimension $r$.
For $r=2$, independent variation in its two factors makes the
two diagonal brackets span both central coordinates; thus
$Z'=Z=\mathbb G_m^2$. The unipotent group is the product
of two Heisenberg groups, with the same pure group acting
on both standard factors.
For $r=1$, after multiplying the points by a common integer
we may write $Q=aP$, $a\in F^\times$. The bracket image is
the one-dimensional line with proportional central coordinates
\begin{equation}\hypertarget{tag-003H}{}\label{ari:eq:cp-positive-centre}
 (1,c),\qquad c=\frac{a a^\dagger}{d}>0,
\end{equation}
where $\dagger$ is the Rosati involution. On an elliptic curve
$aa^\dagger$ is a positive rational number: a square in the
non-CM case and the positive field norm in the CM case.
Construction 3.4(3) adds a further central direction exactly
when the image of the algebraic lift in the quotient torus
is nontorsion. This is the integer $e$ above.
The unipotent Lie algebra is the Heisenberg algebra, with one
additional central Tate line when $e=1$. Formula
\eqref{ari:eq:cp-MTdim} proves every entry of the table. This
description computes the group even when the lift contributes
a Kummer factor; specifying only the span of $P,Q$ would
not determine that last factor.
\end{enumerate}
\end{proof}





% Compendium-only ari:C.7.5: apparatus: verification/diagnostic record retained with source computation




\begin{proposition}[The period-conjecture implication for the local norm]\hypertarget{tag-001D}{}\label{ari:C.7.4}
\StatementRecord{ari:C.7.4}{proved(H)}{001D}{\SeriesRef{ari:C.7.2}, \SeriesRef{ari:C.7.3}, \SeriesRef{ari:eq:cp-v}, \SeriesRef{ari:eq:cp-positive-centre}; \cite[Construction~3.4, p.~12, Lemma~3.5, p.~13, Conjecture~4.7, p.~18, Proposition~4.12 and Remark~4.13, pp.~19--20]{bertolin2022}}
Let $E/\overline\Q$ be an elliptic curve and let
$P\in E(\overline\Q)$ be nontorsion. Use the local norm $N(P)$,
the algebraic extension $G_P$, and its algebraic fibre coordinate
from \cref{ari:C.7.1}, \cref{ari:C.7.2}.
The hypothesis $H$ is the Relative Semi-abelian Conjecture
\cite[Conjecture~4.7, p.~18]{bertolin2022}, restricted to
$G=G_P\times G_{\bar P}$ and the algebraic lifts used below.
Explicitly, for those lifts and their associated $B_Q,Z(1)$, it gives
\[
 \operatorname{trdeg}_{\overline\Q(\Omega_G)}
 \overline\Q\bigl(\Omega_G,R,\bar R,
                  \widetilde{\log}_G(R,\bar R)\bigr)
 \geq 2\dim B_Q+\dim Z(1),
\]
where $\Omega_G$ is the full de Rham--Betti period matrix,
including the second-kind periods. Under $H$, $N(P)$ is
transcendental. This applies to every algebraic parameter
$t\ne0,1$ with nontorsion $P_t$, including the conductor-$361$
CM point identified below. Hypothesis $H$ remains open for
these points. Grothendieck's period conjecture for the algebraic
$1$-motives in this proof implies $H$
\cite[Proposition~4.12, pp.~19--20]{bertolin2022}.
\end{proposition}

\begin{proof}
\begin{enumerate}[label=\arabic*.,leftmargin=*]
\item Suppose that $N(P)$ is algebraic. The compact-character
identity of \cref{ari:C.7.2} supplies an algebraic point
$R=(P,r)\in G_P(\C)^c$: we take the positive algebraic number
$r=|\wp'(u)|\,|\Delta|^{-1/6}N(P)^2$ in the specified
algebraic fibre coordinate. The conjugate point gives the
corresponding algebraic lift $\bar R$ over $\bar P$.
Set $M_R=[\Z\to G_P]$, $1\mapsto R$, and
$\mathcal M=M_R\oplus\overline{M_R}$.

\item Put $K_0=\overline\Q(\Omega_G)$. Write
$P=[u]$ with periods $\omega_1,\omega_2$, quasi-periods
$\eta_1,\eta_2$, and the logarithmic coordinate $v$ of
\eqref{ari:eq:cp-v}. The full matrix contains
\[
 \kappa_j=\zeta(u)\omega_j-\eta_j u,\qquad j=1,2.
\]
The coefficient matrix for $(\zeta(u),u)$ is
$\left(\begin{smallmatrix}\omega_1&-\eta_1\\
\omega_2&-\eta_2\end{smallmatrix}\right)$; its determinant
is the nonzero Legendre determinant. Hence $u,\zeta(u)$
belong to $K_0$. The same calculation in the conjugate factor
places their conjugates there. This is the explicit elliptic
instance of \cite[Remark~4.13, p.~20]{bertolin2022}.

\item Write $u=a\omega_1+b\omega_2$ with $a,b\in\R$.
Solving this system together with its conjugate gives
\begin{equation}\hypertarget{tag-003I}{}\label{ari:eq:cp-real-linear}
 \eta(u)=
 \frac{(u\bar\omega_2-\bar u\omega_2)\eta_1
       +(\bar u\omega_1-u\bar\omega_1)\eta_2}
      {\omega_1\bar\omega_2-\bar\omega_1\omega_2}.
\end{equation}
The denominator is nonzero by real-linear independence of the
periods. Thus $u\eta(u)$ and its conjugate belong to $K_0$.
The compact-character equation is
$\Re(v-u\zeta(u)+u\eta(u))=0$, so
\begin{equation}\hypertarget{tag-003J}{}\label{ari:eq:cp-relation}
 v+\bar v=u\zeta(u)+\bar u\overline{\zeta(u)}
                -u\eta(u)-\bar u\overline{\eta(u)}\in K_0.
\end{equation}
Each conjugate is a coordinate of the conjugate algebraic
$1$-motive, whose full periods were included in $K_0$.

\item The period torsor of $\mathcal M$ maps to the period
torsor of $[0\to G]$. Projection of the abelian point-extension
variety $B$ to the extension-coordinate variety is an isogeny:
the point coordinates are the same algebraic data as the
extension coordinates under the specified polarizations.
Equivalently, in the product $E\times\bar E$, the projected
point $(P,\bar P)$ is the sum of the two extension parameters
$(P,0)$ and $(0,\bar P)$. The formula for $B_Q$ in
\cite[equation~(4.4), p.~19]{bertolin2022} gives $\dim B_Q=0$.
There is therefore no weight-$-1$ kernel. The full kernel on
the connected motivic groups is the weight-$-2$ vector group
$\operatorname{Lie}Z$, as in
\cite[Construction~3.4 and Lemma~3.5, pp.~12--13]{bertolin2022}.
After fixing the base periods its fibre is an affine torsor
over $K_0$ of dimension $\dim Z$.

\item In the positive Tate basis of each algebraic factor,
the central action changes $(v,\bar v)$ by
$2\pi i(c_1,c_2)$. Replacing a conjugated negative Betti loop
by the positive loop changes its period column, while leaving
the analytic coordinate $\bar v$ fixed. The polarization
transport through $\phi:\bar E\to E$ of degree $d$ is
$\phi^tJ_E\phi=dJ_{\bar E}$. Thus the bracket line computed
in \cref{ari:C.7.3} has coordinates
$(1,aa^\dagger/d)$ with $aa^\dagger/d>0$; in the independent
case the bracket fills both central coordinates. In both cases
$(c_1,c_2)\mapsto c_1+c_2$ is nonzero on the actual central
group. Any extra quotient Tate direction preserves this
nonzero restriction. Consequently \eqref{ari:eq:cp-relation}
cuts a proper affine hyperplane in that torsor and forces the
relative transcendence degree to be at most $\dim Z-1$.
The same doubled construction applies to real data; its
rank-one bracket still has positive coefficient.

\item All noncentral generalized logarithm coordinates are in
$K_0$ by steps 2--3, and $R,\bar R$ are algebraic. Hypothesis
$H$ gives the lower bound $\dim Z$, contradicting step 5.
This proves the stated implication. The source's
Proposition~4.12 supplies it from the period conjecture for
these algebraic $1$-motives. No assertion of $H$ is used
elsewhere to remove its stated open status.
\end{enumerate}
\end{proof}

\subsection{CM theta normalization and further unconditional deductions}

The holomorphic theta normalization proposed for the CM interpolation approach differs from the
norm occurring in the stated problem. We make that distinction explicit,
including the finite-character ambiguity in the multiplication law.

\begin{proposition}[the exact CM multiplication law]\hypertarget{tag-001F}{}\label{ari:C.7.6}
\StatementRecord{ari:C.7.6}{proved}{001F}{the definitions and constructions stated in the proof}

Let $E/\overline{\Q}$ have complex multiplication by an order
$\mathcal O$, and use an algebraic invariant differential. Write
$A=\operatorname{area}(\Lambda)$ and
\[
 \eta(z)=s_2z+\frac{\pi}{A}\bar z,\qquad
 \theta(z)=\sigma(z)e^{-s_2z^2/2}.
\]
Then $s_2\in\overline{\Q}$. For every nonzero
$\alpha\in\mathcal O$, with $d=\alpha\bar\alpha$, there is
$F_\alpha\in\overline{\Q}(E)$ such that
\begin{equation}\hypertarget{tag-003Q}{}\label{ari:eq:ct-CMmult}
 \theta(\alpha z)^2=F_\alpha(\exp_E z)\theta(z)^{2d}.
\end{equation}
Its divisor and leading Laurent coefficient are
\[
 \operatorname{div}(F_\alpha)=2[\alpha]^*(O)-2d(O),\qquad
 F_\alpha(\exp_Ez)=\alpha^2z^{2-2d}(1+O(z^2)).
\]
Thus the unsquared quotient is an algebraic function on $E$, but
it need not be a rational function on $E$. For a nontorsion
$P=\exp_Eu\in E(\overline{\Q})$, algebraicity of $\theta(u)$
would imply algebraicity of $\theta(\alpha u)$ for all
$\alpha\in\mathcal O$. The additive group $\mathcal O u$ has
rank two, not more than four.
\end{proposition}


\begin{proof}
\begin{enumerate}[label=\arabic*.,leftmargin=*]
\item The Legendre relation gives the displayed real-linear decomposition
of $\eta$. The sigma automorphy formula gives
\[
 \theta(z+\omega)=\epsilon(\omega)
  \exp\left(\frac{\pi\bar\omega}{A}
                         (z+\omega/2)\right)\theta(z),
 \qquad \epsilon(\omega)\in\{1,-1\}.
\]
For $\alpha\Lambda\subset\Lambda$, the exponential multipliers
of $\theta(\alpha z)/\theta(z)^d$ cancel because
$\alpha\bar\alpha=d$. Its remaining multiplier is a sign.
Squaring therefore gives a meromorphic function on $E$, with the
stated divisor. This divisor is defined over $\overline{\Q}$.
The existence of a rational function with this divisor over $\C$
implies its existence over $\overline{\Q}$: the relevant
one-dimensional Riemann--Roch space is obtained by extension of
scalars. Normalize its leading Laurent coefficient, in the formal
parameter of the algebraic differential, to $\alpha^2$. This
identifies it with the squared analytic quotient. The action of
$\alpha$ on this differential is algebraic, so $\alpha^2$ is an
allowed normalization.

\item Choose a nonreal $\alpha\in\mathcal O$. Since the normalized
sigma series is $\sigma(z)=z+O(z^5)$ with algebraic Laurent
coefficients, the next coefficient of the squared quotient is
\[
 \alpha^2z^{2-2d}
       \{1-s_2(\alpha^2-d)z^2+O(z^4)\}.
\]
The coefficients of the rational function in this algebraic formal
parameter are algebraic. As $\alpha^2-d\ne0$, this proves
$s_2\in\overline{\Q}$ without an unnormalized appeal to a CM
weight-two form. It also proves that the functions in
\eqref{ari:eq:ct-CMmult} are algebraic over the chosen field.
At nontorsion $u$, none of the arguments $\alpha u$ is a lattice
point; evaluating the identity proves the algebraicity implication.
The injective map $\mathcal O\to\C$, $\alpha\mapsto\alpha u$,
proves the rank assertion.

\item The qualification about rationality is necessary. For a square
lattice, $s_2=0$ and $\alpha=1+i$ give
$F_{1+i}=2i\wp$. The unsquared quotient has a simple zero at
the nonzero point in $\ker(1+i)$ and a simple pole at $O$,
which is not a principal divisor on $E$; its square is the
rational function just displayed.
\end{enumerate}
\end{proof}

\begin{cited}[with derivation: two unconditional normalized values]\hypertarget{tag-001G}{}\label{ari:C.7.7}
\StatementRecord{ari:C.7.7}{proved}{001G}{\SeriesRef{ari:C.7.6}, \SeriesRef{ari:C.7.1}, \SeriesRef{ari:eq:ct-Gaussian}, \SeriesRef{ari:C.7.2}, \SeriesRef{ari:C.7.9}; \citet[Theorem~1 and Corollary~1]{waldschmidt1979sigma}; \citet[Theorem~6.1]{huber2022}; \citet[final slide, comment (c)]{bertrand2017}; \citet[p.~15-03]{waldschmidt1988theta}; \citet[Corollary~2]{reyssat1982weierstrass}; \citet[p.~15-05]{waldschmidt1988theta}; \citet[Theorem 3.1.2]{bost2009}; \citet[Theorem 1.4 and Corollary 1.6]{nguyen2023}; \citet[Theorems 1.3 and 2.3]{xie2024algebraicity}}

Let $E/\overline{\Q}$ have an algebraic invariant differential,
and let $u$ be an elliptic logarithm of an algebraic nontorsion
point $P$. Then
\[
 \sigma(u)\exp\left(-\frac12u\zeta(u)\right)
\]
is transcendental. If $E$ has CM and $s_2$ is normalized as in
\ref{ari:C.7.6}, then $\zeta(u)-s_2u$ is also transcendental. The first
assertion concerns the displayed holomorphic normalization;
it does not assert transcendence of either $N(P)$ or
$\theta(u)=\sigma(u)e^{-s_2u^2/2}$.
\end{cited}
The expanded source derivation is retained in the proof compendium, Paper~I, label \nolinkurl{ari:C.7.7}.




The stated normalized values follow from \citet[Theorem~1 and Corollary~1]{waldschmidt1979sigma} and the analytic subgroup theorem in \citet[Theorem~6.1, arXiv v8, p.~43]{huber2022}. A complete re-derivation is given in \CompendiumRef{I}{app:rederive:ari-C.7.7}.


\paragraph{The exact dependence of the theta expression on its lifts.}
Put
$g(u,v)=\log(\sigma(u+v)/(\sigma(u)\sigma(v)))$ on a branch where
all three sigma values are nonzero. Replacing $v$ by $v+\omega$,
$\omega\in\Lambda$, changes $g(u,v)-s_2uv$, modulo $2\pi i\Z$, by
\[
 (\eta(\omega)-s_2\omega)u=\frac{\pi\bar\omega u}{A}.
\]
The map $v\mapsto\exp_Ev$ sends these two lifts to the same
elliptic point. Thus $N\exp_Ev=f(\exp_Eu)$ for a
Rosati-antisymmetric $f$ retains its value under this lift change,
while the displayed theta expression acquires the stated term.
A use of the pointwise criterion in
\citet[final slide, comment~(c)]{bertrand2017} must therefore retain
its compatible lift convention. The Gaussian-norm question here
remains the precise question stated above.

A subsequent source check resolves the normalization of the
one-variable holomorphic assertion itself:
\citet[p.~15-03]{waldschmidt1988theta} explicitly states that
$\sigma(u)e^{-s_2u^2/2}$ is transcendental when the invariants
are algebraic, $E$ has CM, and $\exp_Eu$ is an algebraic
nontorsion point.  That source defines $s_2$ by the regularized
lattice sum and attributes the result to Bertrand's 1982
paper \emph{Valeurs de fonctions th\^eta et hauteurs p-adiques}.
We have checked this complete normalized statement in
Waldschmidt's published exposition; the full original 1982
proof has not been retrieved.  This source clarification
does not validate the unrestricted-lift inference withdrawn
above, and the theorem concerns the holomorphic value rather
than the metric norm.  \ref{ari:C.7.7} retains the two assertions whose
proofs have been fully checked here.

\paragraph{The archimedean term still requiring control.}
The exact comparison is
\begin{equation}\hypertarget{tag-003R}{}\label{ari:eq:ct-Gaussian}
 N(P)=|\Delta|^{1/12}|\theta(u)|
                         \exp\left(-\frac{\pi|u|^2}{2A}\right).
\end{equation}
It follows from \ref{ari:C.7.1} and the real-linear decomposition of
$\eta$. The normalization proved transcendental in \ref{ari:C.7.7} instead
gives the exact identity
\[
 N(P)=|\Delta|^{1/12}
 \left|\sigma(u)e^{-u\zeta(u)/2}\right|
 \exp\left(\frac12\Re\{u(\zeta(u)-\eta(u))\}\right).
\]
Neither formula excludes algebraicity of $N(P)$: transcendence
of one complex factor does not control its absolute value or
its product with the remaining exponential.
The generic finite-order interpolation approach
supplies only the rank-two set $\mathcal O u$; it does not
meet the proposed rank-greater-than-four criterion. The
derivative approach would need the algebraicity of
$\theta'(u)/\theta(u)=\zeta(u)-s_2u$, which \ref{ari:C.7.7} excludes.
Replacing the norm by the holomorphic theta value would
therefore change the problem precisely by the factor in
\eqref{ari:eq:ct-Gaussian}.

The related independence theorem of
\citet[Corollary~2]{reyssat1982weierstrass} says that two of
$e^{i\pi u/\omega_1}$, $s(u)$ and $s'(u)$ are algebraically
independent, where $s(z)=\sigma(z)e^{-\eta_1z^2/(2\omega_1)}$.
For real $u$, hypothetical algebraicity of $N(P)$ makes
$s(u)$ algebraic, so the conclusion instead forces independence
of $e^{i\pi u/\omega_1}$ and
$\zeta(u)-\eta_1u/\omega_1$; it yields no contradiction.
The periodic-functional-equation theorem reported in
\citet[p.~15-05]{waldschmidt1988theta} assumes algebraic
$q=e^{i\pi\tau}$, whereas this $q$ is transcendental for
imaginary quadratic $\tau$ by Gelfond--Schneider.  These
specific checked applications therefore do not decide the
metric norm; their source scopes and the larger CM
independence tuples are recorded separately.

We also tested an additive-extension route.
\citet[Theorem 3.1.2]{bost2009} proves that an algebraic point
of the universal vector extension of an abelian variety which
lies in its maximal compact subgroup at one complex embedding
is torsion. The theorem's hypotheses require an algebraic
point of the vector extension, equivalently an algebraic line
bundle with an algebraic connection having unitary monodromy.
\ref{ari:C.7.2} supplies only an algebraic point in one fibre of a
multiplicative extension. It does not supply that algebraic
connection. Differentiating a metric value at one algebraic
point cannot establish algebraicity of the connection there
or globally. This is the exact missing implication in that
attempt.

Another route uses the Latt\`es map $f_E$ defined by
$x(2Q)=f_E(x(Q))$. Its canonical height is
$\widehat h_{f_E}(x(P))=2\widehat h_O(P)$. For rational
parameters C.6 therefore makes algebraicity of $N(P)$
equivalent to algebraicity of
$\exp(\widehat h_{f_E}(x(P)))$.
\citet[Theorem 1.4 and Corollary 1.6]{nguyen2023} settle the
corresponding multiplicative-height problem for nonintegrable
polynomials using B\"ottcher coordinates in the basin of
infinity. They do not apply to $f_E$. Indeed an elliptic
duplication Latt\`es map has no exceptional point: the inverse
images of any point contain the images of all successive
division points of an elliptic point, an infinite set. A
polynomial has the exceptional point $\infty$, so no conjugate
of $f_E$ is polynomial. The Julia set is the whole sphere
and the required basin of infinity is absent. This checks
the failed hypothesis, rather than treating a theorem for
polynomials as a theorem for all rational maps.

We also checked the arithmetic algebraicity criterion of
\citet[Theorems 1.3 and 2.3]{xie2024algebraicity}.
Besides sufficiently rapid height-normalized approximation by
points in a fixed number field, it requires the analytic germ
to be formally defined over that field. This hypothesis fails
for the compact leaf in \ref{ari:C.7.9}: its tangent has the proved
transcendental slope
$c=\zeta(u)-\eta_1u/\omega_1$.
A smooth formal germ defined over $\overline{\Q}$ has an
algebraic tangent; algebraic translations, coordinate changes
and isogenies preserve this obstruction. Assuming $N(P)$
algebraic supplies an algebraic point on the leaf and its
integer multiples, but not an algebraic formal germ. Attempting
to force convergence by successive division also loses the
fixed-number-field hypothesis: if $[2]^nQ_n=Q$ with all
$Q_n\in E(K)$ for one number field, then
$\widehat h_O(Q_n)=4^{-n}\widehat h_O(Q)>0$.
The $Q_n$ are distinct with bounded ample height, contrary to
Northcott finiteness. The full hypothesis audit and the opened
preprint are retained in the records.

\subsection{The precise nonarchimedean comparison}
\begin{cited}[with derivation: a nonzero ordinary height]\hypertarget{tag-001H}{}\label{ari:C.7.8}
\StatementRecord{ari:C.7.8}{proved}{001H}{\SeriesRef{ari:C.7.6}, \SeriesRef{ari:C.7.3}, \SeriesRef{ari:C.7.4}, \SeriesRef{ari:C.7.2}; \citet[Appendix A]{agboola2006}; \citet{agboola2006}; \citet[Theorem 2.1]{bertrand1995}; \citet[Theorem 3.1]{bertrand1995}; \citet{bost2009}}

Let $E/\Q$ have complex multiplication by the full ring of integers
of an imaginary quadratic field $K$ of class number one. Assume
that its functional-equation sign is $-1$, and let $p>3$ be a
prime of good ordinary reduction, $p\mathcal O_K=\mathfrak p
\mathfrak p^*$. For a nontorsion $P\in E(\Q)$ the
$\mathfrak p$-adic height and the cyclotomic $p$-adic height
of $P$ over $K$, in the normalization of
\citet[Appendix A]{agboola2006}, are nonzero.
In particular this applies to
\[
 E:y^2+y=x^3-38x+90,\qquad P=(0,9),\qquad p=5.
\]
\end{cited}
The expanded source derivation is retained in the proof compendium, Paper~I, label \nolinkurl{ari:C.7.8}.




The nonvanishing theorem is reproduced in Rubin's Appendix~A, Theorem~A.1 and Lemma~A.2 of \citet{agboola2006}; the exact hypotheses and application are retained. A complete re-derivation is given in \CompendiumRef{I}{app:rederive:ari-C.7.8}.


The comparison explains the missing step in an attempted
archimedean adaptation. The ordinary $p$-adic sigma construction
uses a holomorphic quadratic correction. In the CM normalization
the corresponding complex holomorphic function is $\theta$
of \ref{ari:C.7.6}, whereas the real metric also contains
$\exp(-\pi|u|^2/(2A))$. Theorem~A.1 supplies no identity
between the ordinary $p$-adic height and that real factor.
Transferring the conclusion would require precisely such a
new argument; replacing the real factor by the holomorphic
CM correction does not preserve the norm being asked about.

There is also a definite limitation on the maximal-compact
approach. \citet[Theorem 2.1]{bertrand1995} constructs
nontrivial multiplicative extensions of polarized abelian
varieties with nontorsion algebraic points in their maximal
real compact subgroup. Its hypotheses include an endomorphism
not fixed by the Rosati involution and a nontorsion rational
point over a real number field. The construction uses a base
point related to the extension class by a Rosati-antisymmetric
endomorphism. It cannot be replaced by a theorem saying that
every algebraic compact point of every multiplicative
extension is torsion. For our diagonal class, an equality
$NP=f(P)$ with $N>0$ and $f+f^\dagger=0$ would force
$P$ torsion: $N-f$ has no zero eigenvalue, since
$(N-f)^\dagger(N-f)=N^2-f^2$ is positive. In the elliptic
case this is simply a nonzero isogeny. The diagonal positivity
used in \ref{ari:C.7.3}--\ref{ari:C.7.4} is therefore essential, but these
counterexamples neither prove nor refute the desired
unconditional diagonal statement. The additive result in
\citet[Theorem 3.1]{bertrand1995}, and its higher-dimensional
extension in \citet{bost2009}, requires a point in a vector
extension, a hypothesis not furnished by \ref{ari:C.7.2}.


\subsection{The fixed semiabelian extension and the diagonal subgroup attempt}

\begin{proposition}[the endomorphism and slope obstructions]\hypertarget{tag-001I}{}\label{ari:C.7.9}
\StatementRecord{ari:C.7.9}{proved}{001I}{\SeriesRef{ari:C.7.1}, \SeriesRef{ari:eq:cd-compatibility}, \SeriesRef{ari:C.6.R1}, \SeriesRef{ari:C.7.2}; \citep[Theorem 6.1]{huber2022}}

Let $E/\overline\Q$ be an elliptic curve, let
$q\in E^\vee(\overline\Q)$ be nontorsion, and let
\[
 1\longrightarrow\mathbb G_m\longrightarrow G_q
       \longrightarrow E\longrightarrow0
\]
be the extension of class $q$. Then
$\operatorname{End}_{\overline\Q}(G_q)=\Z$, including when
$E$ has complex multiplication. In particular a noninteger
CM endomorphism of $E$ does not lift to an endomorphism of this
fixed extension.

For
\[
 E:y^2+y=x^3-38x+90,\qquad P=(0,9),\qquad Q=2P=(4,-2),
\]
one has
\[
 N(Q)=19^{1/4}N(P)^4.
\]
Use the differential $dx/(2y+1)$, choose the positive real
period $\omega_1$, and choose the real logarithm
$u\in(0,\omega_1/2)$ of $Q$. Put
\begin{equation}\hypertarget{tag-0033}{}\label{ari:eq:cd-slope}
 c=\zeta(u)-\frac{\eta_1}{\omega_1}u,
 \qquad \eta_1=\eta(\omega_1).
\end{equation}
Then $c$ is a real transcendental number. The actual point
\[
 R_N=(Q,3\,19^{-1/2}N(Q)^2)\in G_Q(\C)
\]
is nontorsion, belongs to the maximal compact subgroup, and has
the real logarithm $(u,cu)$. Its algebraicity is equivalent to
the algebraicity of $N(P)$. In the algebraic logarithm coordinates
of \ref{ari:C.7.1}, the line $\ell_c=\C(1,c)$ has transcendental
slope. The proof establishes this slope statement from a period
of $G_Q$ without deciding the algebraicity of $R_N$.
\end{proposition}


\begin{proof}
\begin{enumerate}[label=\arabic*.,leftmargin=*]
\item The torus is the maximal connected affine subgroup of $G_q$.
Consequently every endomorphism of $G_q$ induces multiplication
by some $n\in\Z$ on $\mathbb G_m$ and an endomorphism
$\alpha$ on $E$. Functoriality of extension classes gives
\begin{equation}\hypertarget{tag-0032}{}\label{ari:eq:cd-compatibility}
 \alpha^*q=nq\quad\text{in }E^\vee.
\end{equation}
Here $\alpha^*$ is the dual endomorphism, not $\alpha$
acting on an arbitrarily identified copy of $E$. To make the
convention explicit, choose a principal polarization
$\lambda:E\to E^\vee$, put $Q_0=\lambda^{-1}q$, and
write $\alpha^\dagger=\lambda^{-1}\alpha^*\lambda$.
Equation \eqref{ari:eq:cd-compatibility} says
$(\alpha^\dagger-[n])Q_0=O$. A nonzero endomorphism of an
elliptic curve is an isogeny and has finite kernel. Since
$Q_0$ is nontorsion, $\alpha^\dagger=[n]$, hence
$\alpha=[n]$ as well.

\item The difference between the given endomorphism and $[n]$ has
image in the torus and is trivial on the torus. It therefore
factors through an algebraic-group homomorphism
$E\to\mathbb G_m$. This homomorphism is zero, since a
morphism from the proper connected curve $E$ to the affine
group $\mathbb G_m$ is constant and a group homomorphism
preserves the identity. Thus there is no additional ambiguity
from $\operatorname{Hom}(E,\mathbb G_m)$, and the
endomorphism is exactly $[n]$. Conversely all integer
multiplications exist. This proves the first assertion.

\item For the concrete curve, the tangent slope at $P$ is $-2$,
giving $2P=(4,-2)$. Its discriminant is $-19^3$ and
$\psi_2(P)=2\cdot9+1=19$. The division identity in C.4
therefore gives the asserted formula for $N(Q)$, so replacing
$P$ by $Q$ preserves algebraicity of the norm. This doubling
also places the base point in the identity component of the
real elliptic group without any component assumption. In the
chosen differential the Weierstrass coordinates are
$(\wp,\wp')=(x,2y+1)$, with $g_2=152$, $g_3=-361$.
The negative discriminant means that the real cubic has one
real root. Along $(0,\omega_1/2)$, $\wp$ decreases from
$+\infty$ to that root and $\wp'<0$. Since
$\wp'(u)=2(-2)+1=-3$, the indicated real logarithm of $Q$
exists.

\item The period vectors of $G_Q$ in \ref{ari:C.7.1} include
\[
 (\omega_1,\zeta(u)\omega_1-\eta_1u)
       =\omega_1(1,c).
\]
These are algebraic logarithm coordinates: $dz$ and the
invariant differential $d\log r-df_u/f_u$ are algebraic,
because $df_u/f_u$ is the algebraic third-kind differential
displayed in \ref{ari:C.7.1}. Real invariants and real $u,\omega_1$
make $c$ real.

\item Suppose $c$ were algebraic. Then $\ell_c$ would be an
algebraic Lie subspace. Dividing the nonzero period
$\omega_1(1,c)$ by a sufficiently large positive integer
gives a vector in $\ell_c$ whose exponential is a nonzero
torsion point of $G_Q$, hence an algebraic point. The
analytic subgroup theorem
\citep[Theorem~6.1, arXiv v8, p.~43]{huber2022} supplies a positive-dimensional
connected algebraic subgroup with Lie algebra contained in
$\ell_c$. Its dimension is one, and its projection to $E$
is nontrivial. Such a subgroup must be an elliptic curve
$E'$ projecting isogenously to $E$: a connected
one-dimensional affine algebraic group cannot map
nontrivially to an elliptic curve. Its inclusion splits the
pullback of $G_Q$ to $E'$. The dual of that isogeny would
therefore kill the extension class $\lambda(Q)$, forcing
$Q$ torsion. This contradicts \ref{ari:C.6.R1}. We conclude that $c$
is transcendental. In particular the line is not defined over
$\overline\Q$. This argument uses the period vector alone;
it does not assume algebraicity of $N(P)$.

\item Define the positive real number
\[
 r=3\,19^{-1/2}N(Q)^2.
\]
Proposition~\ref{ari:C.7.2} places the actual point $R_N=(Q,r)$
in the maximal compact subgroup of $G_Q$. Since $Q$ is algebraic,
$R_N$ is algebraic exactly when $r$ is algebraic, equivalently
when $N(Q)$, and hence $N(P)$, is algebraic. The chosen real
logarithm has
\[
 f_u(u)=-\wp'(u)\sigma(u)^2e^{-u\zeta(u)}
        =3\sigma(u)^2e^{-u\zeta(u)}>0.
\]
Consequently $v=\log(r/f_u(u))$ is real and
$(u,v)$ is a logarithm of $R_N$, with no imaginary logarithm
ambiguity. Substituting this real logarithm into the exact
compactness equation $\chi=0$ in \ref{ari:C.7.2} gives
\[
 v=\zeta(u)u-u\eta(u)
   =u\left(\zeta(u)-\frac{\eta_1}{\omega_1}u\right)=cu.
\]
Its projection $Q$ is nontorsion, so $R_N$ is nontorsion.
The input missing from this proposed use of the analytic
subgroup theorem is a proper Lie subspace defined over
$\overline\Q$; the actual compact slope line is proved
not to have that property. The CM decomposition makes the
same obstruction visible as
\[
 c=\zeta(u)-s_2u-\frac{\pi}{A}u.
\]
The last term is the derivative of the real Gaussian
correction. It cannot be discarded while retaining the
compactness equation.
\end{enumerate}
\end{proof}

\paragraph{What functional transcendence gives in this construction.}
The functional theorem can be applied fully, but its conclusion
does not contradict this hypothetical algebraic compact point.
Indeed consider
\[
 f:\C\longrightarrow\operatorname{Lie}(G_Q),\qquad
 f(z)=(z,cz).
\]
No proper connected algebraic subgroup of $G_Q$ projects
onto $E$: as in the proof above, such a subgroup would split
the nontorsion extension after an isogeny. Every proper
connected subgroup is therefore contained in the torus,
whereas $\exp_{G_Q}f$ has nonconstant elliptic projection.
Thus $f$ is $G_Q$-nondegenerate in Definition 1.2 of
\citet{noguchi2024}. Their Theorem 1.3 gives
\[
 \dim_{\C}\overline{\{(\exp_{G_Q}f(z),f(z)):z\in\C\}}^{\rm Zar}
       \ge \dim G_Q+1=3.
\]
The graph is contained in $G_Q\times\ell_c$, of dimension
three, so equality holds. This is a statement over $\C$
about a family of meromorphic functions. It remains valid
whether or not the single value $\exp_{G_Q}f(u)$ is
algebraic. Under the norm-algebraicity hypothesis the integer
values $\exp_{G_Q}f(nu)=[n]R$ are algebraic, but they do
not give an algebraic relation holding on the entire graph.
There is no contradiction with its full dimension.
One would need an additional arithmetic specialization
argument, or a way to obtain a further functional relation
from that single metric value. Neither is supplied by the
checked functional theorem. The endomorphism calculation
above also prevents obtaining extra CM multiplication
directions inside this fixed $G_Q$ by lifting the
noninteger endomorphisms of $E$.


\subsection{A CM-stable enlargement and its Gaussian eigencoordinate}

The fixed-extension calculation in \ref{ari:C.7.9} does not prevent
enlarging the torus. We test that alternative explicitly.
Retain its curve, its point $Q=2P$, and its real logarithm
$u$. Write $\beta=\sqrt{-19}=i\sqrt{19}$, $d=19$,
$D=\zeta(u)-s_2u$, and $k=\pi/A$. All endomorphisms and
extensions in what follows are defined over $\overline\Q$.

\begin{proposition}[the enlarged CM action does not isolate the Gaussian]\hypertarget{tag-0018}{}\label{ari:C.7.10}
\StatementRecord{ari:C.7.10}{proved}{0018}{\SeriesRef{ari:C.7.9}, \SeriesRef{ari:C.7.6}, \SeriesRef{ari:eq:ce-zeta}}

Let
\[
 \mathcal G=G_Q\mathbin{\times_E}G_{\beta Q},\qquad
 1\longrightarrow\mathbb G_m^2\longrightarrow\mathcal G
       \longrightarrow E\longrightarrow0.
\]
There is an endomorphism $B$ of $\mathcal G$ inducing
$\beta$ on $E$ and
$(t_1,t_2)\mapsto(t_2^{-1},t_1^d)$ on the torus.
It satisfies $B^2=[-d]$. Moreover
\[
 \operatorname{Hom}(\mathcal G,\mathbb G_m)=0,
 \qquad
 \operatorname{End}(\mathcal G)\otimes\Q=\Q(\beta).
\]
Thus $\mathcal G$ has no algebraic $\mathbb G_m$ quotient,
including after isogeny.

There is an algebraic compact point in the second extension
above $Q$ whose logarithm is
\[
 (u,v_2),\qquad v_2=-\beta uD+au,
 \qquad
 a=\frac1{2\beta}
       \frac{d}{du}\log F_\beta(Q)\in\overline\Q.
\]
Here the derivative means differentiation of the rational
function $F_\beta(\exp_E z)$ with respect to the algebraic
invariant differential, evaluated at $z=u$.
Pair this point with the actual compact point $R_N$ of
\ref{ari:C.7.9}. The resulting compact point of $\mathcal G(\C)$
has logarithm
\[
 U=(u,v_1,v_2),\qquad v_1=uD-ku^2.
\]
Its algebraicity is equivalent to the algebraicity of $N(P)$.
In these algebraic logarithm coordinates the differential of
$B$ and its $-\beta$ eigenprojection are
\[
 dB=\begin{pmatrix}
 \beta&0&0\\ a&0&-1\\ \beta a&d&0
 \end{pmatrix},\qquad
 \Pi_-:=\frac12\left(I-\frac{dB}{\beta}\right),
\]
\[
 \Pi_-U=-\frac{\pi u^2}{2A}(0,1,\beta).
\]
The first torus coordinate of $\exp(\Pi_-U)$ is exactly
the Gaussian $\exp(-\pi u^2/(2A))$. The projector is not
the differential of an algebraic quasi-endomorphism. No
sequence of rational polynomials in $B$, even allowing
inverses of isogenies, can converge to it as complex linear
operators. At least one of the two torus coordinates of
$\exp(\Pi_-U)$ is transcendental; the proof does not decide the
algebraicity of its first coordinate alone.
\end{proposition}


\begin{proof}
\begin{enumerate}[label=\arabic*.,leftmargin=*]
\item Use the principal polarization to identify the two extension
classes with $Q$ and $\beta Q$. They are linearly independent
over $\Q$ modulo torsion: any nonzero $m+n\beta$ is an
isogeny, so $(m+n\beta)Q$ cannot be torsion. A character
of $\mathbb G_m^2$ given by $(m,n)\in\Z^2$ extends to
$\mathcal G$ exactly when the pushout extension has zero
class, namely when $mQ+n\beta Q=O$. Hence the restriction
of every homomorphism $\mathcal G\to\mathbb G_m$ to
the torus is zero. Such a homomorphism then factors through
$E$ and is zero. This proves the asserted absence of
characters. An isogenous semiabelian variety cannot acquire
a nonzero character, because composing it with the isogeny
would give a nonzero character of $\mathcal G$.

\item The Rosati adjoint of $\beta$ is $-\beta$. Pullback by
$\beta$ sends the ordered extension classes to
$(-\beta Q,dQ)$, exactly as does pushout by the displayed
torus matrix. Equality of these extension classes provides
the endomorphism $B$. It is unique because
$\operatorname{Hom}(E,\mathbb G_m^2)=0$. Its square
induces $[-d]$ on the torus and on $E$, so the same
uniqueness gives $B^2=[-d]$. The map
\[
 \operatorname{End}(\mathcal G)\otimes\Q
     \longrightarrow\operatorname{End}(E)\otimes\Q
          =\Q(\beta)
\]
is injective: an endomorphism inducing zero on $E$ has image
in the torus, and the absence of characters makes it zero.
Its image contains $\beta$, hence equals $\Q(\beta)$.

\item Next construct the compact point without assuming anything
about $N(P)$. The identities in \ref{ari:C.7.6}, applied at $\beta$
and $1+\beta$, give the algebraic nonzero number
\begin{equation}\hypertarget{tag-0036}{}\label{ari:eq:ce-ribet-r}
 r_2=\frac{\theta((1+\beta)u)}{\theta(u)\theta(\beta u)},
 \qquad r_2^2=\frac{F_{1+\beta}(Q)}{F_\beta(Q)}.
\end{equation}
Indeed $(1+\beta)(1+\bar\beta)=1+d$, so the powers
of $\theta(u)$ cancel. None of the arguments is a lattice
point since $Q$ is nontorsion. Written in sigma functions,
\[
 r_2=\frac{\sigma((1+\beta)u)}{\sigma(u)\sigma(\beta u)}
                         e^{-s_2\beta u^2}.
\]
The point $(Q,r_2)$ in the rational trivialization of
$G_{\beta Q}$ therefore has logarithm
\[
 (u,v_2),\qquad
 v_2=u\{\zeta(\beta u)-s_2\beta u\}.
\]
Differentiating \ref{ari:C.7.6} shows that
\begin{equation}\hypertarget{tag-0037}{}\label{ari:eq:ce-zeta}
 \zeta(\beta u)-s_2\beta u=-\beta D+a,
 \qquad a\in\overline\Q,
\end{equation}
with the value of $a$ in the statement. The compactness
character of this extension vanishes, since
\[
 \Re\{v_2-\zeta(\beta u)u+\beta u\eta(u)\}
    =\Re\left(\frac{\beta\pi|u|^2}{A}\right)=0.
\]
This also gives a direct, fully normalized construction of
the antisymmetric compact lift; no unrestricted theta
criterion is being used.

\item For completeness, the matrix $dB$ can be recovered without
choosing a transcendental change of coordinates. Its torus
block is the stated integer matrix and its first entry is
$\beta$, so initially write its first column as
$(\beta,a_1,a_2)^t$, with $a_1,a_2$ algebraic.
The identity $B^2=[-d]$ gives $a_2=\beta a_1$.
For $\omega\in\Lambda$, set
\[
 \kappa_1(\omega)=\zeta(u)\omega-\eta(\omega)u,
 \quad
 \kappa_2(\omega)=\zeta(\beta u)\omega-\eta(\omega)\beta u.
\]
Equation \eqref{ari:eq:ce-zeta} and the real-linear formula for
$\eta$ imply
\[
 \kappa_1(\beta\omega)=a\omega-\kappa_2(\omega).
\]
Comparing the second coordinates of the corresponding period
vectors under $B$ gives
$(a_1-a)\omega_j\in2\pi i\Z$ for $j=1,2$.
Since $\omega_2/\omega_1$ is nonreal, both integers must
be zero, so $a_1=a$. This proves the matrix formula.

\item The actual first compact logarithm in \ref{ari:C.7.9} is
$v_1=uD-ku^2$. Pairing the two compact lifts over the same base
point gives $U$ with no algebraicity assumption. The second lift
is algebraic, so the paired point is algebraic exactly when its
first lift is algebraic.
Substitution into the displayed matrix yields
\[
 v_2+\beta v_1-au=-\beta ku^2,
 \qquad
 \Pi_-U=-\frac{ku^2}{2}(0,1,\beta),
\]
which proves the exact Gaussian identity.

\item The two eigenspaces of $dB$ have eigenvalues $\beta$ and
$-\beta$. Every algebraic quasi-endomorphism has differential
$a_0I+b_0dB$ for rational $a_0,b_0$, with eigenvalues
$a_0+b_0\beta$ and $a_0-b_0\beta$. The projector
$\Pi_-$ would require these values to be $0$ and $1$,
which is impossible. There is also a uniform obstruction to
approximation. For arbitrary real $a_0,b_0$, put
$x=a_0+b_0\beta$, so that its other eigenvalue is
$\bar x$. Then
\begin{equation}\hypertarget{tag-0035}{}\label{ari:eq:ce-gap}
 \max\{|x|,|\bar x-1|\}\ge\frac12,
\end{equation}
because $1\le|\bar x|+|\bar x-1|=|x|+|\bar x-1|$.
These are exactly the two eigenvalue errors in approximating
$\Pi_-$. The spectral radius is bounded by every induced
operator norm, so no such sequence can converge to the
projector. Compositions and inverses of isogenies still belong
to the same endomorphism algebra and satisfy this bound.

\item Finally the image of $\Pi_-$ is the toric Lie line
$\C(0,1,\beta)$. It is not the Lie algebra of an algebraic
subtorus: an algebraic one-dimensional subtorus of
$\mathbb G_m^2$ has rational slope. Since $u\ne0$,
the projected vector is nonzero. The analytic subgroup
theorem in fact shows that its exponential cannot be an
algebraic point of both torus coordinates. Indeed its nonzero logarithm lies in the algebraic line
$\C(1,\beta)$, whereas a one-dimensional algebraic subtorus
has a rational slope. The analytic subgroup theorem would put
a positive-dimensional subtorus in that line for an algebraic
exponential, contradicting these two slope computations. The construction isolates the Gaussian in
a Lie eigencoordinate, without producing an algebraic
group map that evaluates it.

\item The spectral comparison with the ordinary prime $5$ is
explicit. The CM number $\alpha=(1+\beta)/2$ satisfies
$\alpha^2-\alpha+5=0$. Its two archimedean conjugates
both have absolute value $\sqrt5$. At the two places over
$5$, their normalized valuations are $0$ and $1$, because
the polynomial reduces to $X(X-1)$ and has product $5$.
The separation available in that nonarchimedean linear
algebra is absent in \eqref{ari:eq:ce-gap}. This comparison
does not turn a limiting operation on complex Lie spaces
into an algebraic-group endomorphism or transfer a
$p$-adic height conclusion to the real Gaussian.

\item \textbf{An exact compact witness on the conductor-$361$ curve.}
In the preceding normalization the constants and the second
fibre coordinate are explicitly
\begin{equation}\hypertarget{tag-0034}{}\label{ari:eq:ce-explicit-witness}
 s_2=-2,\qquad a=-\frac{21\beta}{703},\qquad
 r_2=-\frac{165}{74}+\frac{1839\beta}{1406}.
\end{equation}
We give polynomial certificates for these values, so they
do not depend on recognition from numerical approximations.
Put $f=4x^3-152x+361$ and
\[
 \begin{aligned}
 h={}&x^9-38x^8+437x^7-1444x^6-7942x^5\\
     &+82308x^4-274360x^3+390963x^2-130321x-130321,\\
 N={}&(-19x+76)h^2+f(h''h-(h')^2)+(6x^2-76)h'h,\\
 K={}&N'h-2Nh'.
 \end{aligned}
\]
Direct expansion gives
\[
 -K^2f=4N^3-152\cdot19^2Nh^4+361\cdot19^3h^6,
 \qquad \gcd(N,h)=1.
\]
Hence, in coordinates $Y=2y+1$, the map
\[
 (x,Y)\longmapsto
 \left(\frac{N}{19h^2},\frac{KY}{19\beta h^3}\right)
\]
is a degree-$19$ endomorphism fixing $O$, with pullback of
$dx/Y$ equal to $\beta\,dx/Y$. It is therefore the CM
endomorphism $\beta$. Divisors and the leading coefficient
at $O$ give $F_\beta=\beta^2h^2$. Its Laurent coefficient
of relative degree two is $-76$. Comparison with \ref{ari:C.7.6}
gives $38s_2=-76$, proving $s_2=-2$.

\item For the other kernel put $h_+=(H_0+\beta H_1)/10$, where
\[
 \begin{aligned}
 H_0={}&10x^8-380x^7+4864x^6-28880x^5+72561x^4\\
       &+54872x^3-795644x^2+1824494x-1433531,\\
 H_1={}&20x^7-304x^6+760x^5+10469x^4-83752x^3\\
       &+246924x^2-315514x+130321.
 \end{aligned}
\]
Writing $\bar h_+$ for replacement of $\beta$ by $-\beta$,
the finite polynomial identities and coprimality checks are
\[
 \begin{gathered}
 19xh^2-N=5f h_+\bar h_+,\qquad
 h_+\mid K+19\beta h^3,\\
 \gcd(h_+,fh)=\gcd(h_+,h_+')=1.
 \end{gathered}
\]
At a root of $h_+$ the displayed CM map sends $(x,Y)$ to
$(x,-Y)$, so these are precisely the eight pairs of
non-$2$-torsion points in $\ker(1+\beta)$. The other three
nonzero kernel points are $E[2]\setminus\{O\}$, since
$1+\beta=2\alpha$ in the CM order. Its norm is $20$,
which verifies the count. Thus
\[
 F_{1+\beta}=\frac{(1+\beta)^2}{4}f h_+^2.
\]
Taking the square ratio in \eqref{ari:eq:ce-ribet-r} and fixing
the sign by its expansion at $u=0$ gives the exact rational
function
\[
 r_2=-\frac{1+\beta}{2\beta}\frac{\wp'(u)h_+(\wp(u))}
                                             {h(\wp(u))}.
\]
At $Q$, the needed evaluations are
\[
 h(4)=-37,\qquad h'(4)=7,\qquad
 h_+(4)=\frac{829-279\beta}{10},\qquad\wp'(u)=-3.
\]
They give exactly \eqref{ari:eq:ce-explicit-witness}. In
particular $(Q,r_2)$ is an explicit nontorsion algebraic
point in the maximal compact subgroup of $G_{\beta Q}$.
Its torus-coordinate norm is
$r_2\bar r_2=974799/26011=4-\wp(\beta u)$.

\item The exact polynomial identities, the complete kernels, the
invariant differential and all displayed algebraic evaluations are
reproduced by
\path{code/current_audits/audit_paper_i/arithmetic_06/cm_witness_exact.py}.
Its JSON output records $39$ exact checks, including the two good-reduction
point counts used to exclude rational torsion. For analytic corroboration,
put
\[
 \tau=\frac{1+i/\sqrt{19}}2,\qquad
 \omega_1=\left(\frac{g_2(\tau)}{152}\right)^{1/4}>0,
 \qquad u=\omega_1\wp_\tau^{-1}(4\omega_1^2),
\]
using the branch $0<u<\omega_1/2$ with $\wp'(u)=-3$.
The real-period basis is obtained from
$i(1,\tau_0)$, $\tau_0=(1+i\sqrt{19})/2$, by the integral
basis change $(1,\tau_0)\mapsto(1-2\tau_0,1-\tau_0)$, whose
matrix has determinant one; its period ratio is the displayed $\tau$.
Evaluation with interval arithmetic at $160$ decimal working digits in
\path{code/current_audits/audit_paper_i/arithmetic_06/cm_analytic_balls.py}
gives
\[
 \begin{aligned}
 c&=1.14793408195385447403276278660663356664777041472753\ldots,\\
 e^{-\pi u^2/(2A)}
  &=0.68133647525194864004469591167703191986051023921164\ldots.
 \end{aligned}
\]
Its output records $10$ checks. In particular the differences between the
sigma expression for $r_2$ and its exact value, and between the zeta
expression for $a$ and its exact value, each have absolute value less than
$10^{-110}$. The polynomial identities prove the algebraic values; this
numerical comparison corroborates the stated analytic branch.
\end{enumerate}
\end{proof}

\subsection{Complexifying the Gaussian: values, jets, and the rank criterion}

The antiholomorphic Gaussian can be made holomorphic by introducing a
second variable. This gives more algebraic values under the hypothesis
in C.7(c), but the derivative and growth hypotheses of the resulting
transcendence arguments must still be checked. We do this on the
conductor-$361$ model, where all constants and conjugations are fixed.

\begin{proposition}[the complexified route and its exact obstructions]\hypertarget{tag-0019}{}\label{ari:C.7.11}
\StatementRecord{ari:C.7.11}{proved}{0019}{\SeriesRef{ari:C.7.9}, \SeriesRef{ari:C.7.1}, \SeriesRef{ari:C.7.6}, \SeriesRef{ari:eq:cc-fixedfield}, \SeriesRef{ari:C.7.7}; \citep[Theorem 6.1]{huber2022}}

Use the curve, point and real logarithm of \ref{ari:C.7.9}:
\[
 E:y^2+y=x^3-38x+90,\qquad Q=(4,-2)=\exp_Eu,
 \qquad 0<u<\omega_1/2.
\]
Put $\beta=i\sqrt{19}$, $k=\pi/A$, $s_2=-2$, and
\[
 \theta(z)=\sigma(z)e^{-s_2z^2/2},\qquad
 D(z)=\zeta(z)-s_2z,\qquad
 \mathcal F(z,w)=\sqrt{19}\,\theta(z)\theta(w)e^{-kzw},
 \qquad H=\mathcal F^2.
\]
Then $\mathcal F(z,\bar z)=N(z)^2$, and $\mathcal F$ is an
entire function of two complex variables of order at most two.
It is periodic under $(z,w)\mapsto(z+\omega,w+\bar\omega)$,
$\omega\in\Lambda$. The subgroup
\[
 \Gamma=\{(\alpha u+\omega,\bar\alpha u+\bar\omega):
                   \alpha\in\Z[\beta],\ \omega\in\Lambda\}
 \subset\C^2
\]
has rank four over $\Z$ and lies in the real plane
$\{(z,\bar z):z\in\C\}$. There is one number field $K$ such
that, with $H_0=H(u,u)=N(Q)^4$, every value of $H$ on
$\Gamma$ belongs to $K[H_0]$. More precisely, the nonzero values
are algebraic multiples of the powers $H_0^{\alpha\bar\alpha}$,
as specified in \eqref{ari:eq:cc-fixedfield}.

For every nonzero $\alpha\in\Z[\beta]$ there is an algebraic
number $a_\alpha$ such that
\[
 D(\alpha u)=\bar\alpha D(u)+a_\alpha.
\]
At $p_\alpha=(\alpha u,\bar\alpha u)$ the algebraic constant
direction $V_\alpha=\alpha\partial_z-\bar\alpha\partial_w$
satisfies
\begin{align*}
 V_\alpha\log\mathcal F(p_\alpha)
    &=\alpha a_\alpha-\bar\alpha\bar a_\alpha
                         \in\overline{\Q},\\
 V_\alpha^2\log\mathcal F(p_\alpha)
    &=-\alpha^2\{\wp(\alpha u)+s_2\}
      -\bar\alpha^2\{\wp(\bar\alpha u)+s_2\}
      +2\alpha\bar\alpha k.
\end{align*}
The number $k$ is transcendental. For every nonzero $\alpha$,
the exact normalized second derivative satisfies
\[
 \overline\Q\!\left(\frac{V_\alpha^2H(p_\alpha)}{H(p_\alpha)}\right)
      =\overline\Q(k).
\] In particular,
\[
 (\partial_z-\partial_w)H(u,u)=0,\qquad
 \frac{(\partial_z-\partial_w)^2H(u,u)}{H(u,u)}=4k-8.
\]
Thus the normalized first derivative is algebraic and the
normalized second derivative generates $\overline\Q(k)$. Moreover the rank-four set does not satisfy the
explicit derivative-free finite-order criterion audited below.
These are obstructions to the stated applications of those
criteria; they do not decide the algebraicity of $N(Q)$.
\end{proposition}



\begin{proof}
\begin{enumerate}[label=\arabic*.,leftmargin=*]
\item The real model and its algebraic differential give real sigma
coefficients and a conjugation-stable lattice. In particular
$\overline{\theta(z)}=\theta(\bar z)$. The discriminant is
$-19^3$, so \ref{ari:C.7.1} gives the asserted diagonal identity.
The sigma function and the exponential of the quadratic
polynomial $-kzw$ have order at most two. From \ref{ari:C.7.6},
\[
 \theta(z+\omega)=\epsilon(\omega)
         e^{k\bar\omega(z+\omega/2)}\theta(z),
 \qquad \epsilon(\bar\omega)=\epsilon(\omega)\in\{1,-1\}.
\]
Multiplying this formula and its conjugate version, the two
signs multiply to one and the exponent cancels the change of
$-kzw$. This proves periodicity on all of $\C^2$, not only
on the real diagonal.

\item If $\alpha u+\omega=0$, then $[\alpha]Q=O$.
Nontorsion of $Q$ and the fact that a nonzero elliptic
endomorphism is an isogeny imply $\alpha=0$, and then
$\omega=0$. Thus the rank-two module $\Z[\beta]u$ and the
rank-two period lattice are independent over $\Z$.
This proves the rank assertion, with no extra algebraic
independence assumption on $u$ and the periods.

\item For $d=\alpha\bar\alpha$, the squared CM law \ref{ari:C.7.6} gives
the exact fixed-field identity
\begin{equation}\hypertarget{tag-002Q}{}\label{ari:eq:cc-fixedfield}
 H(\alpha u,\bar\alpha u)
   =19^{1-d}F_\alpha(Q)\overline{F_\alpha(Q)}H(u,u)^d.
\end{equation}
All the endomorphisms and their normalized rational functions
$F_\alpha$ are defined over a single number field: choose
a field containing the coefficients of the finitely many
generators of the CM order, and close it under conjugation.
The uniqueness of the Laurent normalization in \ref{ari:C.7.6} gives
the same field of definition for every $F_\alpha$.
Writing that fixed number field as $K$, formula
\eqref{ari:eq:cc-fixedfield} places all values in $K[H_0]$,
with no algebraicity assumption on $H_0$.
The value for $\alpha=0$ is zero, and period shifts do not
change it. Squaring was included because algebraicity of the
individual positive square roots alone would not establish
their membership in one fixed number field.

\item Logarithmically differentiating the squared CM law at $u$
gives
\[
 a_\alpha=\frac{1}{2\alpha}
     \frac{(d/dz)F_\alpha(\exp_Ez)|_{z=u}}{F_\alpha(Q)}
                  \in\overline{\Q}.
\]
The denominator is nonzero since $Q$ is nontorsion. The
derivation $d/dz$ is the algebraic invariant derivation on
$E$, so the displayed derivative has algebraic value.
Since $D(u)$ is real, conjugation gives the corresponding
formula at $\bar\alpha u$. On a neighborhood where
$\mathcal F\ne0$, its logarithmic gradient and Hessian are
\label{layout:I:3}
\[
\begin{gathered}
(\partial_z\log\mathcal F,\partial_w\log\mathcal F)
       =(D(z)-kw,D(w)-kz),\qquad
 \\
\operatorname{Hess}\log\mathcal F=
 \begin{pmatrix}-\wp(z)-s_2&-k\\-k&-\wp(w)-s_2\end{pmatrix}.
\end{gathered}
\]
Evaluation on $(\alpha,-\bar\alpha)$ proves both identities
in the statement. In particular the Gaussian cancels in the
first derivative and reappears in the second derivative.

\item To prove $k$ transcendental, recall the real period formula
$\eta_1=(s_2+k)\omega_1$. If $k$ were algebraic, the
nonzero period vector $(\omega_1,\eta_1)$ of the algebraic
universal vector extension of $E$ would lie on the algebraic
line of slope $s_2+k$. Dividing this period by a sufficiently
large integer supplies a nonzero algebraic torsion exponential.
The analytic subgroup theorem
\citep[Theorem~6.1, arXiv v8, p.~43]{huber2022} would give a one-dimensional
algebraic subgroup with that Lie line. Its projection to $E$
is nonzero, so it is an elliptic curve isogenous to $E$ and
splits the vector extension after that isogeny. This is
impossible: its nonzero class in $H^1(E,\mathcal O_E)$ stays
nonzero under an isogeny, whose pullback has a left inverse
after dividing the trace by the degree. These are the same
algebraic logarithm coordinates and extension as in \ref{ari:C.7.7};
a fixed algebraic linear change of coordinates leaves the
argument unchanged.

\item Write the two logarithmic derivatives as
$V_\alpha\log\mathcal F=B_\alpha$ and
$V_\alpha^2\log\mathcal F=A_\alpha+2dk$, where
$A_\alpha,B_\alpha$ are algebraic. Since $H=\mathcal F^2$,
\[
 \frac{V_\alpha^2H(p_\alpha)}{H(p_\alpha)}
       =4B_\alpha^2+2A_\alpha+4dk.
\]
Here $d>0$, so the displayed affine expression both belongs to
$\overline\Q(k)$ and recovers $k$ by subtraction of its algebraic
constant and division by $4d$. This proves the asserted field
equality; $H(p_\alpha)$ is nonzero because $Q$ is nontorsion.
At $\alpha=1$, $a_1=0$, $\wp(u)=4$ and $s_2=-2$,
which give the last displayed formula in the statement.

\item \textbf{The precise several-variable criteria that were checked.}
The differential criterion of
\citet[Theorems 3.6 and 3.7]{waldschmidt2009auxiliary}
requires a finite algebraically independent family whose
partial derivatives belong to the polynomial ring generated
by that family over one number field. Algebraic values of
the full family then give algebraic values of all iterated
partial derivatives. The calculation above gives the exact value-field obstruction:
the field containing $H(p_\alpha)$ and its second directional
derivative contains $k$ as soon as the nonzero value is inverted.
In particular its normalized second jet is transcendental. The issue persists
although the first cancelling direction is algebraic; the
direction also varies with $\alpha$.

\item For a criterion without differential equations we opened
\citet[Theorem 1, p.~24]{waldschmidt1975dimension}, including
its proof and a visual check of the displayed inequalities.
The source uses $d$ meromorphic functions on $\C^n$, growth and
arithmetic-size exponents $\rho_1,\ldots,\rho_d$, and sampling
sets of cardinality $N^{\ell-o(1)}$ in a ball of radius $O(N)$,
with separation at least $N^{-\lambda-o(1)}$. Its strict numerical
requirement is
\begin{equation}\hypertarget{tag-002X}{}\label{ari:eq:cc-rankcriterion}
 (d-1)\ell>\sum_{j=1}^d\rho_j
                     +2d(n-1)(\lambda+1).
\end{equation}
The theorem also requires all function values to belong to
one number field and bounds on denominators in the entire
representations of the meromorphic functions. We do not
silently discard these hypotheses.

\item The natural independent triple here is
$\wp(z),\wp(w),H(z,w)$, all of order at most two.
The first two are algebraically independent. If $H$ satisfied
a polynomial over $\C(\wp(z),\wp(w))$, fix a generic pair
$(z,w)$ and translate only $z$ by $m\omega$, with
$0\ne\omega\in\Lambda$. The coefficients of the polynomial
remain constant, whereas the absolute values of its proposed
roots $H(z+m\omega,w)$ grow as
$\exp(k|\omega|^2m^2+O(|m|))$. This contradicts the finite
set of roots of a nonzero polynomial. Thus the triple is
independent.

\item The growth-order lower envelope is $\rho_j=2$. Substitution of
this envelope and $n=2$, $d=3$, $\ell=4$ into
\eqref{ari:eq:cc-rankcriterion} gives the required inequality
\[
 8>12+6\lambda,
\]
which fails even for $\lambda=0$. The actual geometry is
less favorable: the $\asymp N^4$ points obtained by bounding
the four integer coefficients lie in a ball of radius $O(N)$
in a real plane. Packing separated discs in that plane gives
$N^4\ll N^{2+2\lambda+o(1)}$, hence $\lambda\ge1$.
Removing the rank-two set of poles does not change this count.
Larger growth exponents only worsen the inequality. This
checks the particular published criterion rather than asserting
that every possible multivariable transcendence method fails.
The missing unconditional step in C.7(c) remains the exclusion
of an algebraic value on the diagonal compact leaf.
\end{enumerate}
\end{proof}

\subsection{Division values and a variable-degree interpolation test}


\begin{proposition}[division ratios, scalar models and translated functions]\hypertarget{tag-001A}{}\label{ari:C.7.12}
\StatementRecord{ari:C.7.12}{proved}{001A}{\SeriesRef{ari:C.4.2}, \SeriesRef{ari:C.7.1}, \SeriesRef{ari:C.7.6}, \SeriesRef{ari:C.7.9}, \SeriesRef{ari:C.7.11}}

Let $E:y^2+y=x^3-38x+90$, $Q=(4,-2)$ and $N_0=N(Q)$.
Use the positive real logarithm $u$ and period $\omega_1$ of
\ref{ari:C.7.9}. Put
\[
 \delta=19^{1/4},\qquad b=\eta_1/\omega_1,\qquad
 g(z)=\delta\sigma(z)e^{-bz^2/2}.
\]
Thus $|g(z)|=N(\exp_Ez)$ on the real axis. Let $\mathcal D$
be the CM division hull of $Q$. For $R\in\mathcal D$, the
equations $[n]R=[\alpha]Q$, $n>0$, determine a unique
$\kappa_R=\alpha/n\in\operatorname{End}(E)\otimes\Q$.
Set $w_R=\kappa_R\bar\kappa_R$. For $R\ne O$,
\[
 A_R:=\frac{N(R)}{N_0^{w_R}}\in\overline\Q_{>0}.
\]
In particular, for $[n]R=[m]Q$ with $m\ne0$,
\begin{equation}\hypertarget{tag-003S}{}\label{ari:eq:division-ratio-model}
 A_R^{n^2}=\delta^{n^2-m^2}
                  \frac{|\psi_m(Q)|}{|\psi_n(R)|}.
\end{equation}
For each positive algebraic number $\lambda$, define
\[
 N_\lambda(R)=A_R\lambda^{w_R}
       =N(R)\left(\frac{\lambda}{N_0}\right)^{w_R}.
\]
All these values are algebraic. For $[n]R=[m]Q$, $|m|\le n$,
$R\ne O$, there are constants $D,H$ depending on $\lambda,E,Q$
such that
\[
 h(N_\lambda(R))\le H,\qquad
 [\Q(R,N_\lambda(R)):\Q]\le
 \begin{cases}Dn^4,&R\in E(\R),\\2Dn^6,&R\in E(\C).
 \end{cases}
\]
The transformation is exact for every positive real parameter:
$N_\lambda(R)/N_\mu(R)=(\lambda/\mu)^{w_R}$ and
$N_{N_0}(R)=N(R)$.

For any finite collection of real division-hull logarithms $a$,
the functions $g(z)$ and $g(z+a)/g(a)$, with $\exp_Ea\ne O$,
have function-field transcendence degree at most three over
$\overline\Q$ after adjoining the elliptic coordinates.
The three functions $g(z),g(z+v)/g(v),g(z-v)/g(v)$, $v=u/q$,
attain three. The normalized holomorphic extensions
$F_a(z)/N(a)^2$ of $N(\exp_E(z+a))^2$ for fixed complex
division-hull shifts belong to the same type of field.
The unnormalized translates have the same upper bound, attained by
the displayed triple, over the algebraic closure $L$ of $\overline\Q(N_0)$ in $\C$.

For the scalar models on the real division grids
$z=(mu+k\omega_1)/n$, $|m|\le Cn$, $k$ modulo $n$, with
fixed $C$, every fixed finite family of these values has a common
degree envelope $O(n^5)$ and bounded heights. For the actual values the same construction
gives a relative degree envelope over a fixed finite extension
of $\Q(N_0)$. The algebraicity of $N_0$ remains the open question
of \ref{ari:C.7.4}. The explicit interpolation and tangent
calculations below record what these bounds establish.
\end{proposition}

\begin{proof}\leavevmode\par


 
\begin{enumerate}[label=\arabic*.,leftmargin=0pt,labelsep=.5em]
\item \emph{The algebraic ratios and their exact change of parameter.}
Two equations $[n]R=[\alpha]Q$ and $[n']R=[\alpha']Q$
give $[n'\alpha-n\alpha']Q=O$. A nonzero elliptic endomorphism
has finite kernel, whereas $Q$ is nontorsion. Thus
$\alpha/n=\alpha'/n'$, proving uniqueness. The resulting
map $R\mapsto\kappa_R$ is additive, and its kernel is the
torsion subgroup. Every such $R$ is algebraic, since it lies
in a finite isogeny fibre above an algebraic point.


\begin{enumerate}[label={$\langle2\rangle$\arabic*.},leftmargin=1.3em,labelsep=.5em]
\item  The division identity and the squared CM law give, for
$d=\alpha\bar\alpha$ and $\alpha\ne0$,
\[
 N(R)^{n^2}=\frac{N(nR)\delta^{n^2-1}}{|\psi_n(R)|},
 \qquad
 N(\alpha Q)^2=\delta^{2-2d}|F_\alpha(Q)|N_0^{2d}.
\]
Consequently
\[
 A_R^{2n^2}=\frac{\delta^{2n^2-2d}|F_\alpha(Q)|}
                         {|\psi_n(R)|^2}\in\overline\Q_{>0}.
\]
The positive root proves $A_R\in\overline\Q$. For integral
$\alpha=m$, the usual multiplication identity is
$N(mQ)=N_0^{m^2}|\psi_m(Q)|/\delta^{m^2-1}$;
substitution proves \eqref{ari:eq:division-ratio-model}.
For nonzero $n$-torsion $R$, take the limit at the simple zero
of $\psi_n$ in the uniformizing coordinate. Since
$N(z)=\delta|z|+o(|z|)$ at zero,
\[
 A_R^{n^2}=N(R)^{n^2}
       =\frac{n\delta^{n^2}}{|\partial_z\psi_n(R)|}.
\]
The invariant derivation is algebraic, so this also is algebraic.
All exponents $w_R$ are rational. The asserted scalar models,
their parameter-change identity and the specialization at $N_0$
follow directly. In particular their definition is independent
of every presentation of $R$ by a division equation.


\item  Put $K_\lambda=\Q(\delta,\lambda)$ and
$D=[K_\lambda:\Q]$. For integral $m\ne0$,
\[
 N_\lambda(mQ)=\lambda^{m^2}
                     |\psi_m(Q)|/\delta^{m^2-1}\in K_\lambda,
\]
because $Q$ is rational. The fibre of $[n]$ has $n^2$ points,
so adjoining $R$ costs at most $n^2$. For real $R$ the absolute
value of $\psi_n(R)$ differs from it by a sign. The division
equation for $N_\lambda(R)$ costs at most another $n^2$.
For complex $R$, adjoin $R$ and $\bar R$, at cost at most
$n^4$, and use
\[
 N_\lambda(R)^{2n^2}
 =\frac{N_\lambda(mQ)^2\delta^{2(n^2-1)}}
              {\psi_n(R)\psi_n(\bar R)}.
\]
This gives $2Dn^6$. At nonzero torsion points the derivative
formula proves the same bounds.

 
\end{enumerate}

\item \emph{Uniform heights.}
We give the coefficient estimate needed here, so the height
assertion does not rely on an unspecified division-polynomial
bound.  On the fixed integral model, $\psi_n$ is a polynomial
in $x$ for odd $n$, and $(2y+1)$ times a polynomial in $x$
for even $n$.  These polynomials have degree $O(n^2)$.
Choose a fixed circle $|x|=A>1$ avoiding the finite branch values.
On its compact inverse image in $E(\C)$, $N$ has a positive
minimum $m_A$, and $|2y+1|$ has a positive minimum.  Let $M$
be the maximum of $N$ on the whole compact elliptic curve.
The rearranged division identity gives
\[
 |\psi_n(R)|\leq M\delta^{n^2-1}m_A^{-n^2}
       \qquad(|x(R)|=A).
\]
Cauchy's coefficient estimate, and division by $2y+1$ for
even $n$, show that the logarithm of the coefficient
$\ell^1$ norm is $O_E(n^2)$.  The same assertion holds for
$\partial_z\psi_n$, because differentiation and the fixed
Weierstrass equation add only polynomial factors in $n$.


\begin{enumerate}[label={$\langle2\rangle$\arabic*.},leftmargin=1.3em,labelsep=.5em]
\item  The elementary height estimate for a fixed degree-four
duplication map, followed by telescoping, gives
$h(x(R))=2\widehat h(R)+O_E(1)$ uniformly over
$E(\overline\Q)$.  Since
$\widehat h(R)=(m/n)^2\widehat h(Q)$, the coordinate heights
are uniformly bounded when $|m|\leq n$.
At each nonarchimedean place the ultrametric inequality bounds a
polynomial value by its largest monomial; at each archimedean
place use the sum of the monomials. Summing these logarithmic
bounds by the product formula gives the polynomial evaluation
height inequality, and hence
$h(\psi_n(R))=O_{E,Q}(n^2)$, and the same estimate for its
invariant derivative at torsion points.  Also
\[
 h(N_\lambda(mQ))\leq m^2h(\lambda)+O_{E,Q}(m^2).
\]
Apply the same division identity to $N_\lambda$ and divide its height by $n^2$.
For a complex value one uses
$h(|\alpha|)\leq h(\alpha)$, which follows by squaring and
using $|\alpha|^2=\alpha\bar\alpha$.
This proves a uniform bound for $h(N_\lambda(R))$.

 
\end{enumerate}

\item \emph{Normalized translates and their algebraic scalars.}
For a real logarithm $a$ of a nonzero division point put
\[
 h_a(z)=\frac{\sigma(z+a)}{\sigma(z)\sigma(a)}e^{-baz}.
\]
The exact identities are
\[
 g(z+a)=\delta^{-1}g(z)g(a)h_a(z),\qquad
 h_a(z+\omega_1)=h_a(z),\quad
 h_a(z+\omega_2)=e^{-2\pi ia/\omega_1}h_a(z).
\]
They follow by substitution from sigma automorphy and Legendre's
relation. A real point in the CM division hull lies in the
ordinary division hull: conjugating $[n]R=[\alpha]Q$ gives
$[\alpha-\bar\alpha]Q=O$, hence $\alpha=\bar\alpha\in\Z$.
The real group is connected because the discriminant is negative.
Its division logarithms are therefore exactly
$\widetilde{\mathcal D}=\Q u+\Q\omega_1$.
The coefficient map
\[
 \nu:\widetilde{\mathcal D}\longrightarrow\Q,
 \qquad \nu(ru+s\omega_1)=r
\]
is well defined and additive: a rational relation between $u$
and $\omega_1$ would make $Q$ torsion.


\begin{enumerate}[label={$\langle2\rangle$\arabic*.},leftmargin=1.3em,labelsep=.5em]
\item  For any positive real $\lambda$ define, on this additive group,
\[
 \widetilde g_\lambda(z)
   =g(z)(\lambda/N_0)^{\nu(z)^2},\qquad
 h_{a,\lambda}(z)
   =h_a(z)(\lambda/N_0)^{2\nu(z)\nu(a)}.
\]
At nonzero division values, the absolute value of
$\widetilde g_\lambda$ is $N_\lambda$ and its sign is the real
sign of $g$. Thus these values are algebraic for algebraic
$\lambda$. The identity
\[
 h_{a,\lambda}(z)=
 \frac{\delta\widetilde g_\lambda(z+a)}
      {\widetilde g_\lambda(z)\widetilde g_\lambda(a)}
\]
is proved by expanding $(\nu(z)+\nu(a))^2$.
It specifies the exact transformation on the sampling group.


\item  Choose a common denominator $q$ for finitely many shifts,
write $a_j=(m_ju+k_j\omega_1)/q$, and put $v=u/q$.
The quotient $h_{a_j}/h_v^{m_j}$ is elliptic for
$\Lambda_q=\Z\omega_1+\Z q\omega_2$. This finite isogeny
of $E$ has a model over $\overline\Q$ with the stated
differential. The quotient has an algebraic divisor, hence
is a rational function over $\overline\Q$ times a scalar.
The scalar is algebraic. Indeed the parameter factors cancel:
\[
 \frac{h_{a_j}(z)}{h_v(z)^{m_j}}
  =\frac{h_{a_j,1}(z)}{h_{v,1}(z)^{m_j}}
       \qquad(z\in\widetilde{\mathcal D}).
\]
Choose a real division point avoiding the finite zeros and poles.
The right side and the value of the chosen rational function
are algebraic and nonzero, so their quotient is algebraic.
This proves that $g(z)$ and all normalized real translates are
algebraic over $\overline\Q(E)(g,h_v)$, after the fixed finite
isogeny. The upper transcendence-degree bound is three.


\item  For the lower bound, work even over $\C$. The multiplier of
$h_v$ under $q\omega_2$ is $e^{-2\pi iu/\omega_1}$,
which has infinite order. A polynomial relation over the elliptic
field, evaluated at a generic $z$ and translated by $mq\omega_2$,
would have infinitely many distinct roots obtained by these
powers. Thus $h_v$ is transcendental over that field. Moreover
\[
 |g(z+mq\omega_2)|=\exp(c m^2+O(|m|)),
 \qquad c=\pi q^2\Im(\omega_2/\omega_1)>0.
\]
In a proposed polynomial relation for $g$ over the elliptic field
and $h_v$, clear denominators. The nonzero leading coefficient
is a polynomial evaluated along an irrational rotation; it is
bounded away from zero on a subsequence, and all other
coefficients remain bounded there. Quadratic growth makes the
highest power of $g$ dominate, a contradiction. Finally
\[
 \wp(z)=\wp(v)-
 \frac{\delta^2g(z+v)g(z-v)}{g(z)^2g(v)^2}
\]
recovers the elliptic coordinate from the three stated translates,
and $h_v=\delta g(z+v)/(g(z)g(v))$. They therefore attain three.
The constants $g(a)$ lie in $L$, by the already proved ratios.
Multiplication by these constants gives the exact map from
normalized to raw translates, proving the assertion over $L$.


\item  For a complex CM division shift $a$, the holomorphic extension
from the real axis is
\[
 F_a(z)=\frac{N(a)^2}{\delta^2}g(z)^2h_a(z)h_{\bar a}(z).
\]
Its multiplier depends on $a+\bar a\in\Q u+\Q\omega_1$.
Choose $q,m$ with $a+\bar a=(mu+k\omega_1)/q$.
Then $h_ah_{\bar a}/h_v^m$ is elliptic on $\C/\Lambda_q$
with algebraic divisor. Its scalar is algebraic by the same
evaluation argument, using the exact identity on real division
points
\[
 h_a(z)h_{\bar a}(z)(\lambda/N_0)^{2\nu(z)m/q}
 =\frac{\delta^2N_\lambda(\exp_E(z+a))^2}
        {N_\lambda(\exp_Ea)^2\widetilde g_\lambda(z)^2}.
\]
Here $w_{\exp_E(z+a)}=|\kappa_{\exp_Ea}+\nu(z)|^2$;
expanding the square proves the identity. Set $\lambda=1$
and divide by $h_{v,1}^m$. This proves the normalized complex
translate assertion with all scalar factors accounted for.

 
\end{enumerate}

\item \emph{Joint degrees and heights on the sampling group.}
In the algebraic meromorphic section coordinate of
\ref{ari:C.7.1},
\[
 z\longmapsto(\exp_Ez,h_v(z))
   =\exp_{G_V}\bigl(z,(\zeta(v)-bv)z\bigr),\qquad
 V=\exp_Ev.
\]
For the scalar model, multiply its torus coordinate by
$(\lambda/N_0)^{2\nu(z)\nu(v)}$. This factor is a character
of the additive sampling group, so the resulting map
$\Phi_\lambda:\widetilde{\mathcal D}\to G_V(\C)$ is an
abstract group homomorphism. Its real period is $\omega_1$,
because $\nu(\omega_1)=0$ and the displayed analytic map
has that period. For algebraic $\lambda$, its values in the displayed section
chart are algebraic by the formula for $h_{v,\lambda}$ above.
For an exceptional sampling point at a zero or pole of that
section, write $z=z_1+z_2$ with both $z_j$ avoiding the finite
exceptional set modulo the elliptic lattice. Such a choice exists
in the infinite division group. Then
$\Phi_\lambda(z)=\Phi_\lambda(z_1)\Phi_\lambda(z_2)$ is
algebraic by the algebraic group law. This also covers the identity.


\begin{enumerate}[label={$\langle2\rangle$\arabic*.},leftmargin=1.3em,labelsep=.5em]
\item  For $z=(mu+k\omega_1)/n$,
$[n]\Phi_\lambda(z)=[m]\Phi_\lambda(u)$.
Choose a fixed number field containing the coefficients of
$G_V$, the point $\Phi_\lambda(u)$ and the finitely many
fixed constants. Multiplication by $n$ on $G_V$ has degree
$n^3$, since its kernel is an extension of $E[n]$ by
$\mathbb G_m[n]$. Adjoining $\widetilde g_\lambda(z)$ costs
at most $n^2$ by the norm division equation. Passing to the
fixed isogeny and adjoining fixed constants changes only the
constant in $O(n^5)$. The same argument for any positive real
$\lambda$ works over a fixed finite extension of $\Q(\lambda)$,
because all the fixed division values are algebraic over that
field. At $\lambda=N_0$ it gives the claimed relative envelope.


\item  The height proof above applies with $|m|\le Cn$ for any fixed
$C$, changing only its constant. It bounds the model values at
$z$ and every fixed real translate $z+a$. Their quotients bound
$h_{v,\lambda}(z)$. Coordinates on the fixed isogenous elliptic
curve have bounded height by the fixed isogeny height comparison.
The normalized complex translates are fixed rational functions
of those coordinates, $\widetilde g_\lambda^2$ and a fixed power
of $h_{v,\lambda}$, with algebraic coefficients just proved.
The polynomial evaluation inequality bounds their heights.
Any removable zero or pole in these expressions lies in a fixed
finite set and is handled by its actual model value. Taking the
positive square root for a fixed complex-shifted norm divides
height by two and costs at most a factor two in degree.

 
\end{enumerate}

\item \emph{The exact numerical and tangent checks.}\par
The sufficient inequality in Gramain--Mignotte--Waldschmidt
\citep[Theorem2.1,pp.101--102]{gramain1986values} has the form
\[
 d(N)\prod_{j=1}^k\max\{d(N)\psi_j(N),\theta_j(R(N))\}
 \le c_0\mu(N)^{k-1}\varphi(N)^k,
\]
\[
 \varphi(N)=\log\frac{R(N)^2+r(N)r(N+1)}
                         {R(N)(r(N)+r(N+1))}.
\]
On fixed nested discs the growth and geometric factors are
constant. The scalar-model height envelopes are constant too.
The numerical comparison of these envelopes is therefore
$d^{k+1}\ll\mu^{k-1}$. At denominator $n$ the real group
has $O(n^2)$ sample points, and through denominator $N$ it
has $O(N^3)$. Substitution of the proved envelopes gives
$N^{12}$ against $N^3$ for two functions and $N^{20}$ against
$N^6$ for three. These upper-envelope calculations provide
no such inequality and assert no lower bound for the actual
degrees. For actual analytic samples the base field is the
displayed extension of $\Q(N_0)$; its number-field status is
exactly the unresolved algebraicity issue. Extra normalized
translates remain in the function field of transcendence degree
three proved above.


\begin{enumerate}[label={$\langle2\rangle$\arabic*.},leftmargin=1.3em,labelsep=.5em]
\item  The arithmetic-germ criteria
\citep[Theorem6.1]{herblot2012meromorphic}
\citep[Theorem4.16]{gasbarri2010analytic}
can be compared with the actual logarithmic tangent. At a real
nontorsion division point $z$, the exact formula is
\[
 \frac1{g(z)}\frac{dg}{dx}
       =\frac{\zeta(z)-bz}{\wp'(z)}.
\]
Its numerator is transcendental by the period-line proof of
\ref{ari:C.7.9}, applied to this algebraic nontorsion point;
its denominator is nonzero algebraic. A formal graph defined
over $\overline\Q$ would have an algebraic nonzero constant
value and algebraic first derivative, contradicting their
displayed ratio. The scalar change on the sampling group has
not changed this exact tangent of the analytic graph.


\item  For the Mahler system in
\citep[Theorem1.1 and system(1.2)]{adamczewski2017mahler},
the needed algebraic Taylor coefficients and algebraic regular
argument can be tested directly: the coefficient of $z^3$ in
$g$ is $-\delta b/2$, and $b=s_2+k$ is transcendental by
\ref{ari:C.7.11}. The logarithm $u$ is transcendental too:
an algebraic $u\ne0$ would give an algebraic point $(u,Q)$
on the diagonal Lie line of $\mathbb G_a\times E$; the analytic
subgroup theorem would produce an algebraic one-dimensional
subgroup projecting nontrivially to both factors, which is
impossible since an elliptic curve has no nonconstant map to
$\mathbb G_a$ and $\mathbb G_a$ has no nonzero group map to $E$.
The duplication identity acts by the explicit elliptic
multiplication map. The entropy identity in
\citep[Theorem5.1]{einsiedler2001entropy} uses the sequence
built from the denominators of $x(2^nQ)$ with scale $4^n$.
These precise inputs supply no evaluation of $N_0$ here.


\item  Finally the period-refined rank comparison
\citep[Theorem4.1 and Corollary4.3]{waldschmidt1988periodrank}
\citep[main theorem,p.363]{roy1993periodrank}
has coefficient $\alpha(G)=2\dim G-2d_0(G)-d_1(G)$,
where the last terms count additive and multiplicative
quotients. Both vanish for $G_Q$. An additive homomorphism
vanishes on its torus and then on the proper quotient $E$;
a nonzero torus character extending to $G_Q$ would kill a
nonzero integer multiple of its nontorsion extension class.
Thus $\alpha(G_Q)=4$.


\item  The line $\ell_c$ has period rank one. A general period is
$(\omega,\zeta(u)\omega-\eta(\omega)u+2\pi i j)$.
Writing $\omega=a\omega_1+d\omega_2$, membership in
$\ell_c$ gives $du/\omega_1\in\Z$ by Legendre's relation.
Since $u/\omega_1$ is irrational, $d=j=0$. The numerical
rank bound in the cited theorem is
$(4-1)\dim\ell_c=3$. The period and the actual logarithm
$(u,cu)$ are independent over $\Q$ and span rank two;
the algebraicity of the latter's exponential remains exactly
the point-algebraicity question in \ref{ari:C.7.9}. These
explicit ranks leave that question open.


\end{enumerate}

\end{enumerate}
\end{proof}


\section{Open arithmetic questions for the normalized norm}\hypertarget{tag-003T}{}\label{ari:open:norm}



Without hypothesis $H$, the arithmetic nature of the metric norm remains open at the
nontorsion points of \cref{ari:C.7.4}. An explicit small instance is
\[
 E:y^2=x^3-140x+784,\qquad P=(0,28),\qquad j=-3375,
 \qquad N(P)=2e^{-\widehat h(P)}.
\]
Here $P$ is a primitive generator modulo the rational
$2$-torsion and
\[
 \widehat h(P)=
 0.46736567718147128000752020126284678363562034370331\ldots.
\]
The exact unresolved assertion is that $2e^{-\widehat h(P)}$
is transcendental. The corresponding example of conductor
$361$, $y^2+y=x^3-38x+90$, $P=(0,9)$, has
$N=e^{-\widehat h(P)}$ and no finite local contribution.
The relative-period proof isolates its obstruction: an
unconditional theorem excluding
\eqref{ari:eq:cp-relation} for the positive diagonal central
direction. A theorem about linear $1$-period relations
does not supply that exclusion. The attempted approaches and the exact CM theta comparison are printed in the preceding period and CM-normalization sections;
their source checks and computational records are retained in the audit files.
These arguments do not resolve the stated exclusion.

\clearpage



\clearpage
